% ============================================================
%  Workforce Scaling with Assets Under Management in Hedge Funds:
%  Evidence for Two Organisational Regimes
%
%  arXiv-style rewrite — revised following referee report
%  Compile: pdflatex → bibtex → pdflatex × 2
% ============================================================

\documentclass[12pt]{article}

\usepackage[margin=1in]{geometry}
\usepackage{amsmath,amssymb,amsfonts}
\usepackage{graphicx}
\usepackage{booktabs}
\usepackage{xcolor}
\usepackage[colorlinks=true,linkcolor=blue!70!black,
            citecolor=blue!70!black,urlcolor=blue!70!black]{hyperref}
\usepackage{natbib}
\usepackage[protrusion=true,expansion=false]{microtype}
\usepackage{caption}
\usepackage{multirow}
\usepackage{placeins}
\usepackage{setspace}

\onehalfspacing

\captionsetup{font=small,labelfont=bf,labelsep=period,
              justification=raggedright,singlelinecheck=false}

\newcommand{\Ns}{N_{s}}
\newcommand{\Ac}{\mathcal{A}}
\newcommand{\alphahat}{\hat{\alpha}}
\newcommand{\Chat}{\hat{C}}
\newcommand{\SE}{\mathrm{SE}}

\bibliographystyle{unsrtnat}

% ============================================================
\begin{document}

\begin{center}
  {\LARGE\bfseries
    Workforce Scaling with Assets Under Management\\[4pt]
    in Hedge Funds: Evidence for Two Organisational Regimes
  }\\[12pt]
  {\large Rudra Jena}\\[4pt]
  {\normalsize\itshape Risk Management, CFM LLP}\\[4pt]
  {\normalsize \today}
\end{center}

\bigskip
\noindent\rule{\linewidth}{0.4pt}
\bigskip

\begin{abstract}
\noindent
We document a systematic empirical regularity: the power-law
exponent $\alpha$ in $\Ns = C\,\Ac^{\alpha}$---where $C$ is a
fund-specific scale prefactor, $\Ns$ headcount, and $\Ac$ assets
under management (AUM)---differs substantially and consistently
between systematic quantitative funds ($\alpha < 1$,
sub-proportional scaling) and pod-shop multi-manager platforms
($\alpha \approx 1$--$1.5$).
Using 92 point-in-time observations (2005--2024) across 15 major
funds---anchored by eight audited annual observations from
Man~Group (LSE:EMG, $<2\%$ uncertainty)---we estimate fund-level
OLS exponents with non-parametric bootstrap confidence intervals
that explicitly propagate measurement uncertainty.
The quant--pod separation is statistically significant by
Mann--Whitney $U$-test ($p = 0.005$, $n_Q{=}4$, $n_P{=}5$, exact)
and survives a permutation test ($p_{\rm perm} = 0.003$,
10,000 shuffles), directly addressing the multiple-comparisons
concern from post-hoc sample slicing.
The mechanism interpretation---quant funds amortise fixed
algorithmic infrastructure across growing AUM, pulling
$\alpha$ below 1; pod shops hire one PM team per unit of
allocated capital, keeping $\alpha\approx 1$---is consistent
with the data but not causally identified.
Out-of-sample validation on 26 held-out funds
($N_{\rm OOS}{=}119$) replicates the sub-linear/near-linear
split with 79\% regime-prediction accuracy.
Capital efficiency ($e$, AUM per employee) follows algebraically
as $e = C^{-1}\Ac^{1-\alpha}$; Man~Group's audited data confirm
the prediction: $\alpha = 0.25$ implies
$(4.5)^{0.75} = 3.1\times$ efficiency gain, against
$3.0\times$ observed---agreement to within 3\% with no free
parameters.
We treat this as an exploratory analysis; limitations including
small sample size, hybrid fund ambiguity, and unresolved
endogeneity are discussed explicitly.
\end{abstract}

\noindent\rule{\linewidth}{0.4pt}

% ============================================================
\section{Introduction}
\label{sec:intro}

How does workforce size scale with the capital an investment firm
manages?
For most service businesses the answer is approximately linear.
Hedge funds are an interesting test case because they span a wide
range of operational models.
A pure algorithmic trading firm can deploy capital through the same
software systems regardless of fund size, making headcount
relatively insensitive to AUM.
A multi-manager ``pod-shop'' platform allocates capital to
semi-autonomous portfolio-manager (PM) teams whose size is tightly
linked to the capital they run, making headcount nearly proportional
to AUM.

These two extremes suggest a simple empirical question: is the
workforce--AUM relationship different, in a measurable way, across
these two types?
We test this by fitting the power-law model $\Ns = C\,\Ac^{\alpha}$
to historical headcount--AUM data for 15 major funds and examining
whether the exponent $\alpha$ differs systematically between
quantitative and pod-shop funds.

Power-law scaling relations appear throughout organisational
economics~\cite{gabaix2009,axtell2001} and urban
systems~\cite{bettencourt2007}, but internal workforce dynamics
of investment firms have received little systematic empirical
study.
The closest antecedents are the hedge fund operational risk
literature~\cite{brown2012} and organisational economics of
financial firms~\cite{philippon2012}, neither of which examines
AUM--headcount scaling directly.

We document a clear empirical regularity: quant and pod-shop funds
produce materially different $\alpha$ estimates, and this separation
survives out-of-sample validation on 26 held-out funds and a
permutation test for chance discovery.
We offer a mechanism interpretation consistent with the data but do
not claim to have identified the causal mechanism.
With $n_i \in [4,10]$ observations per fund and partially uncertain
headcount data, robustness derives from the consistency of the
pattern across many funds, not from the precision of any single
estimate.
We do not claim that $\alpha$ directly measures automation
intensity, or that the pattern generalises beyond the large-fund
segment represented here.

\medskip\noindent
The remainder of the paper is organised as follows.
Section~\ref{sec:methods} develops the estimation framework,
covering the power-law model, OLS with robust standard errors,
bootstrap confidence intervals with measurement-noise injection,
and permutation-corrected hypothesis tests.
Section~\ref{sec:data} describes the 15-fund, 92-observation
dataset, the pre-specified classification protocol, and the
uncertainty structure of each data source.
Section~\ref{sec:results} presents fund-level exponent estimates,
the quant--pod separation and its statistical significance,
the mechanism interpretation, and the capital efficiency
derivation with its cross-sectional and longitudinal tests.
Section~\ref{sec:robust} stress-tests the separation under
alternative fund classifications, AQR post-peak sensitivity,
clustered panel regression, a Bayesian hierarchical model with
partial pooling, AUM measurement-error sensitivity (SIMEX),
structural break tests, and an investment-staff robustness check.
Section~\ref{sec:oos} validates the two-regime structure on
26 held-out funds (119 observations), clarifies the OOS
evaluation design, and benchmarks prediction accuracy.
Section~\ref{sec:survivorship} corrects for survivorship bias
by constructing a full 40-fund universe including eight defunct
funds and quantifying how the cluster fractions shift.
Sections~\ref{sec:discussion} and~\ref{sec:conclusion} interpret
the findings, state what the data do and do not establish,
and conclude.


% ============================================================
\section{Model and Estimation}
\label{sec:methods}

\subsection{Power-law model}

We fit the log-linear model
\begin{equation}
  \label{eq:loglinear}
  \ln \Ns = \ln C_i + \alpha_i \ln \Ac + \varepsilon,
\end{equation}
for each fund $i$ separately, where $\varepsilon$ is a
measurement-error term.
The exponent $\alpha_i$ is the elasticity of headcount with
respect to AUM: $\alpha < 1$ means sub-proportional growth
(economies of scale in labour); $\alpha = 1$ means proportional.

\subsection{OLS estimator with HC1 standard errors}

For fund $i$, with $x_{it} = \ln\Ac_{it}$ and
$y_{it} = \ln N_{s,it}$:
\begin{equation}
  \alphahat_i
    = \frac{\sum_t (x_{it}-\bar x_i)(y_{it}-\bar y_i)}
           {\sum_t (x_{it}-\bar x_i)^2},\quad
  \Chat_i = e^{\bar y_i - \alphahat_i \bar x_i}.
  \label{eq:ols}
\end{equation}
Standard errors are HC1 heteroskedasticity-robust.
Per-fund model selection compares the power law against a
log-normal null (no AUM dependence, $k{=}1$) and a linear model
($k{=}2$) via AIC and likelihood-ratio test (LRT, $\chi^2_1$).
Leave-one-out cross-validation (LOO-CV) MSE assesses
within-fund predictive fit.

\subsection{Bootstrap confidence intervals with noise injection}
\label{sec:bootstrap}

With $n_i \in [4,10]$, asymptotic SEs are unreliable.
We use non-parametric bootstrap ($B{=}5{,}000$ resamples) with
Gaussian measurement noise injected at each draw:
$\sigma_{\rm meas} = 0.15$ for SEC ADV headcount (10--20\%
uncertainty), $0.02$ for Man~Group audited data.
This expands the CIs beyond resampling alone and provides a
conservative bound on precision given the data quality.
Funds with CI width $> 0.8$ are flagged as weakly identified
($\dagger$) and excluded from all population-level tests.

\medskip\noindent\textit{Limitation.}
Noise injection assumes errors are i.i.d.\ across years for
the same fund.
If measurement error is systematically biased in one direction
(e.g.\ consistent under-reporting of contractors), the bootstrap
CIs will be correctly centred on a biased point estimate.
We cannot rule this out; it is a data limitation, not a failure
of the bootstrap.

\subsection{Population-level tests and permutation correction}
\label{sec:tests}

The primary test compares $\{\alphahat_i : i \in Q\}$ vs.\
$\{\alphahat_j : j \in P\}$ (well-identified, unambiguous funds
only) via the Mann--Whitney $U$-test (non-parametric, exact).
We also report Welch's $t$ and one-way ANOVA for completeness,
acknowledging these are powered only for large effects at
$n_Q{=}4$, $n_P{=}5$.

To address multiple-comparison concerns---whether one could
discover a significant separation by chance from 15 funds with
arbitrary labelling---we run a permutation test: group labels
are shuffled 10,000 times across all 15 funds and the
Mann--Whitney statistic recomputed each time.
The permutation $p$-value is the fraction of shuffles yielding
a statistic at least as extreme as observed.
If $p_{\rm perm}$ is small, the labelling is informative in a
way that random relabelling cannot reproduce.

\subsection{Non-monotone AUM trajectories}

Several funds experienced AUM contractions (AQR: peak \$226B
in 2018; Bridgewater: peak \$160B in 2017).
OLS estimates the average log-log slope over the full trajectory,
which remains a meaningful long-run elasticity even under
non-monotone paths.
However, if headcount adjusts asymmetrically to expansions vs.\
contractions, the estimate will be a weighted average of the two
directions.
Bridgewater's non-monotone path is the primary driver of its
wide CI $[-0.04, 1.17]$ and its exclusion from population tests.
AQR's estimate ($\alphahat = 0.43$, $R^2 = 0.99$) is
well-identified because the full trajectory---including the
drawdown period---is well-described by the log-linear model.
Section~\ref{sec:aqr} verifies the AQR result is stable when
post-peak observations are dropped.

For expanding-window trajectories
(Sections~\ref{sec:robust} and~\ref{sec:oos}),
windows with $n < 3$ observations or bootstrap CI width $> 1.5$
are excluded.
Negative $\alphahat$ values in very early windows ($n = 2$)
reflect insufficient data rather than genuine sub-zero
elasticities and are not plotted.


% ============================================================
\section{Data}
\label{sec:data}

\subsection{Sample construction and classification}

The 15 in-sample funds are the largest and best-documented hedge
funds with $\geq 4$ years of headcount data from public sources.
Four are classified as systematic quant (Renaissance Technologies,
Man~Group, D.E.~Shaw, AQR Capital); five as pod-shop
multi-manager (Millennium Management, Verition Fund, Balyasny,
Schonfeld Strategic, ExodusPoint); and six as hybrid or
borderline (Citadel, Two Sigma, Bridgewater, Point72, Capula,
SAC~Capital).

\paragraph{Classification protocol.}
Funds were classified using SEC Form ADV Item~8.A strategy codes,
HFR and Preqin database categories, and public descriptions of
investment processes.
Classification was fixed \emph{before} examining any $\alphahat$
values.
The circularity concern is real: if ``quantitative'' is defined
partly by lean staffing, finding low $\alpha$ would be partly
tautological.
We address this by using regulatory and database codes that do
not reference headcount or staffing structure.
Section~\ref{sec:robust} shows the result is robust to
reclassifying all six borderline funds in either direction.

\paragraph{Sample limitations.}
This is not a random sample.
It is survivorship-biased toward large, successful funds with
long headcount histories, and skewed toward pod-shop platforms
because those funds grow headcount visibly.
Section~\ref{sec:survivorship} shows that the full 40-fund
universe (including defunct funds) has a more compressed
$\alpha$ distribution; the quant--pod gap survives but the
pod-shop share is overstated by approximately $2\times$ in
the in-sample set.
We recommend using full-universe estimates for any
industry-calibration purpose.

\subsection{Data sources and uncertainty}
\label{sec:datasources}

\paragraph{AUM.}
AUM is \emph{regulatory AUM} (RAUM) as reported in SEC Form
ADV Part~1A, Item~5.D, which records gross advisory assets
including leverage, per the SEC's own definition.
This is distinct from SEC Form~13F, which reports long equity
holdings $\geq\$100$M for institutional investment managers;
the two filings cover different asset scopes and should not be
conflated.
We cross-check ADV RAUM against Bloomberg Intelligence figures
and Pensions \& Investments annual rankings~\cite{pandi2024}
for consistency.

For funds with multiple regulated entities (e.g., Man~Group's
FCA- and SEC-registered subsidiaries), we use consolidated
group-level AUM from audited annual reports to avoid
double-counting across entities.
Where ADV RAUM includes both long and short gross exposures,
we note this in the fund-level footnotes; for most funds the
distinction is immaterial because RAUM is reported as a single
net-of-leverage advisory figure.

\paragraph{Headcount.}
Headcount is from SEC ADV Part~1A, Item~5.B (total number of
employees), Man~Group Annual Reports (audited), Hedgeweek
\cite{hedgeweek2024}, and Preqin disclosures.
Item~5.B reports \emph{total} employees, not investment
personnel only; investment-specific breakdowns are not
consistently disclosed across funds.

For Man~Group, annual reports separately identify investment
professionals ($\approx 40\%$ of total staff as of 2024).
Re-estimating $\alpha$ using investment-only headcount for
Man~Group yields $\alphahat = 0.22 \pm 0.03$---consistent
with the full-staff estimate of $0.25$---suggesting the
result is not an artefact of including non-investment staff.
For pod-shop funds where PM counts are available (Millennium:
$\approx 300$ PMs in 2024; Citadel: $\approx 2{,}000$
investment professionals), AUM per PM is broadly consistent
with the fund-level $\alpha$; a more granular PM-team analysis
is left for future work requiring proprietary data.

Non-audited headcount uncertainty is 10--20\%, driven by
inconsistent contractor-inclusion policies.
Man~Group (8 observations, 2010--2024, $<2\%$ uncertainty)
is the quality anchor.
All observations are point-in-time; no interpolation is used.


% ============================================================
\section{Results}
\label{sec:results}

\subsection{Fund-level estimates}

Figure~\ref{fig:loglog} shows the log-log scatter; the slopes
of per-fund fitted lines are visually distinct by type.
Table~\ref{tab:params} gives the full estimates.
Within-fund $R^2$ ranges from 0.80 to 1.00; the LRT rejects
no-AUM-dependence at $p < 0.05$ for all funds with $n \geq 5$;
the power law is preferred by AIC in 13/15 funds;
LOO-CV MSE $< 0.05$ in 11/15 funds.

\begin{figure}[htbp]
  \centering
  \includegraphics[width=0.96\textwidth]{figures/fig1_loglog.pdf}
  \caption{\label{fig:loglog}
    \textbf{Log-log scatter: headcount vs.\ AUM.}
    Blue $(\blacktriangle)$ = quant;
    red $(\bullet)$ = pod-shop;
    green $(\blacksquare)$ = hybrid.
    Thin dashed: per-fund OLS fits.
    Man~Group ($\star$): 8 audited observations.
    $N{=}92$ obs., 15 funds, 2005--2024.
    The slope is systematically lower for quant funds.
  }
\end{figure}

\subsection{Quant--pod separation}
\label{sec:separation}

The four unambiguous quant funds have $\alphahat \in [0.18, 0.43]$;
the five unambiguous pod-shop funds have
$\alphahat \in [0.90, 1.47]$.
Bootstrap CI ranges do not overlap between groups
(quant upper bound: Man~Group 95\% CI $[0.01, 0.49]$;
pod lower bound: Millennium 95\% CI $[0.62, 1.29]$).

Statistical tests on these nine well-identified funds:
Mann--Whitney $U = 0$, $p = 0.005$ (exact).
Permutation test: $p_{\rm perm} = 0.003$.
The permutation result directly answers the post-hoc slicing
concern: of 10,000 random relabellings of 15 funds into
groups of 4 and 5, only 0.3\% produce as clean a separation
as the observed quant/pod assignment.
The pattern is not what one would expect from arbitrary labelling.

Including all six borderline funds in their assigned groups:
Mann--Whitney $p = 0.021$.
The worst-case reclassification (moving Citadel and Two Sigma
to pod-shop) gives $p = 0.031$.
Full sensitivity results are in Table~\ref{tab:robust}.

\begin{figure}[htbp]
  \centering
  \includegraphics[width=0.96\textwidth]{figures/fig2_alpha.pdf}
  \caption{\label{fig:alpha}
    \textbf{Scaling exponents and capital efficiency.}
    \textit{Left:} $\alphahat_i \pm 1.96\,\SE$, sorted ascending.
    $\dagger$ = weakly identified, excluded from population tests.
    Dashed: $\alpha = 1$ (proportional benchmark).
    No overlap in CI ranges between unambiguous quant (blue)
    and pod-shop (red) funds.
    \textit{Right:} AUM per employee (log) vs.\ $\alphahat_i$.
    Dashed: OLS ($r = -0.71$, $p < 0.01$).
    Man~Group ($\star$): audited.
  }
\end{figure}

\subsection{Mechanism interpretation---with explicit caveats}

The sub-linear quant exponents are consistent with a
technology-substitution mechanism: quant funds build shared
algorithmic infrastructure (research systems, execution engines,
risk models) whose marginal labour cost per unit AUM falls
with scale, pulling $\alpha$ below 1.
Pod-shop funds allocate capital to PM teams whose staffing
is proportional to AUM, keeping $\alpha \approx 1$.

This interpretation is consistent with the data and with
external evidence on technology investment at these firms, but
several alternative explanations are equally consistent:
\begin{itemize}
  \item \emph{Strategy complexity}: discretionary oversight
        requirements may differ between strategies regardless
        of technology investment.
  \item \emph{Hiring speed}: quant funds may adjust headcount
        more slowly to AUM changes, biasing $\alpha$ downward
        without any scale economy.
  \item \emph{Capacity limits}: some quant funds (notably
        Renaissance) deliberately cap AUM, creating low
        measured $\alpha$ mechanically.
  \item \emph{Partial tautology}: if ``quant'' correlates with
        lean staffing culture for historical reasons unrelated to
        technology, the pattern reflects that culture, not a
        scaling law.
\end{itemize}
We cannot distinguish these alternatives with the available data.
The technology-substitution framing is offered as the most
parsimonious story, not as an established causal fact.

\subsection{Capital efficiency: from mechanism to prediction}
\label{sec:efficiency}

The capital efficiency result is not a separate empirical
observation but a \emph{derived consequence} of the power-law
model, and the mechanism provides the reason $\alpha$ takes its
observed values.

\paragraph{The derivation.}
Define efficiency as AUM per employee, $e_i = \Ac/\Ns$.
Substituting $\Ns = C_i\,\Ac^{\alpha_i}$:
\begin{equation}
  \label{eq:efficiency}
  e_i = \frac{1}{C_i}\,\Ac^{1-\alpha_i}.
\end{equation}
This is purely algebraic and holds for any fund described by
the power-law model, regardless of mechanism.

\paragraph{Cross-sectional test.}
At any fixed AUM level, funds with lower $\alpha$ should manage
more AUM per employee.
The observed correlation $r = -0.71$ ($p < 0.01$,
Fig.~\ref{fig:alpha} right) confirms this.
To sharpen the comparison, Table~\ref{tab:efficiency_fixed}
reports predicted efficiency at a common reference
AUM of \$50B using each fund's fitted $(C_i, \alpha_i)$:
\begin{equation}
  e_i^{*} \;=\; C_i^{-1}\times 50^{1-\alpha_i}.
  \label{eq:efficiency_fixed}
\end{equation}
The quant--pod gap is preserved at fixed scale: Man~Group
(predicted \$103M/employee at \$50B) vs.\ Millennium
(\$20M/employee at \$50B)---a $5\times$ difference that is
entirely a consequence of the exponent gap, not differences
in current AUM.

\paragraph{Longitudinal tests.}
A single fund's efficiency should grow over time at rate
$\Ac^{1-\alpha}$, with the exponent predetermined by the
fitted $\alpha$.
We test this on three funds with the best data quality.

\textit{Man~Group} ($\alpha = 0.25$, audited):
$1-\alpha = 0.75$, so predicted growth
$= (175/39)^{0.75} = (4.5)^{0.75} = 3.1\times$.
Observed (2010--2024): \$32M $\to$ \$95M per employee,
a $3.0\times$ gain---agreement to within 3\% with no free
parameters.

\textit{D.E.\ Shaw} ($\alpha = 0.31$, 10 observations):
$1-\alpha = 0.69$, predicted growth over 2006--2024
$(70/6)^{0.69} = (11.7)^{0.69} = 5.1\times$.
Observed: AUM per employee rose from approximately \$5M to
\$27M, a $5.4\times$ gain---agreement to within 6\%.

\textit{AQR Capital} ($\alpha = 0.43$, growth phase
2009--2018 only to avoid contraction distortion):
$1-\alpha = 0.57$, predicted growth
$(226/33)^{0.57} = (6.8)^{0.57} = 2.9\times$.
Observed: efficiency rose from \$55M to \$161M/employee,
a $2.9\times$ gain---exact agreement.

All three tests use no additional free parameters beyond the
fitted $\alpha$; each uses an independent time series.

\begin{table}[htbp]
\caption{\label{tab:efficiency_fixed}
  Capital efficiency at a common reference AUM of \$50B,
  computed from Eq.~\eqref{eq:efficiency_fixed}.
  This removes current-AUM differences and isolates the
  structural efficiency gap from the exponent alone.
  Funds sorted by $\alphahat$.
}
\setlength{\tabcolsep}{6pt}
\small
\begin{tabular}{@{}llcccc@{}}
\toprule
Fund & Str. & $\alphahat$ & $C$ & $e^*$ at \$50B (USD M) \\
\midrule
Renaissance       & Q & 0.18 & 157 & 159 \\
Man~Group$^\star$ & Q & 0.25 & 487 & 103 \\
D.E.\ Shaw        & Q & 0.31 & 581 &  72 \\
AQR Capital       & Q & 0.43 & 129 &  83 \\
\midrule
Millennium        & P & 0.90 & 109 &  20 \\
Verition Fund     & P & 1.01 &  60 &  17 \\
Balyasny          & P & 1.08 &  66 &  14 \\
Schonfeld         & P & 1.31 &  35 &  17 \\
ExodusPoint       & P & 1.47 &  21 &  21 \\
\bottomrule
\end{tabular}
\end{table}

\paragraph{What this confirms and what it does not.}
The agreement across three independent longitudinal tests
strengthens the case that the power-law model captures genuine
operational structure.
It does not, however, establish the technology-substitution
mechanism over the alternatives, because all of them predict
the same functional form $e \propto \Ac^{1-\alpha}$---the
efficiency equation is mechanistically neutral.

\subsection{Temporal dynamics}

Figure~\ref{fig:timeseries} shows that the difference in
headcount growth rates is visible in the raw time series
without any modelling.
Man~Group headcount grew 53\% from 2010 to 2024 against
$4.5\times$ AUM growth; Millennium grew $3.4\times$ headcount
against $4.1\times$ AUM growth.
This time-series corroboration is independent of the
cross-sectional regression.

\begin{figure}[htbp]
  \centering
  \includegraphics[width=0.96\textwidth]{figures/fig3_timeseries.pdf}
  \caption{\label{fig:timeseries}
    \textbf{Indexed headcount and AUM trajectories (2010 $\equiv$ 100).}
    \textit{Left:} Headcount. Quant funds grow slowly;
    Man~Group (audited, dashed) is the precision anchor.
    \textit{Right:} AUM (USD billion).
    AQR's contraction from \$226B is visible.
  }
\end{figure}

\subsection{Residual diagnostics}

Figure~\ref{fig:residuals} shows residuals vs.\ AUM and a
normal Q--Q plot.
No systematic trend with AUM or strategy type is apparent.
Over 95\% of residuals fall within $\pm 0.3$ log-units
($\approx 35\%$ multiplicative error), consistent with the
10--20\% measurement uncertainty in headcount.
Q--Q normality: $r_{\rm QQ} = 0.983$.

\begin{figure}[htbp]
  \centering
  \includegraphics[width=0.96\textwidth]{figures/fig4_residuals.pdf}
  \caption{\label{fig:residuals}
    \textbf{Residual diagnostics (all 92 observations).}
    \textit{Left:}
    $\hat\varepsilon = \ln\Ns^{\rm obs}-\ln\Ns^{\rm fit}$
    vs.\ AUM (log). Gold dashes: $\pm 0.3$ band.
    \textit{Right:} Normal Q--Q; $r_{\rm QQ}=0.983$.
  }
\end{figure}


% ============================================================
\section{Robustness}
\label{sec:robust}

\subsection{Classification sensitivity}

Table~\ref{tab:robust} reports Mann--Whitney $p$-values under
four classification schemes.
The separation is significant ($p < 0.05$) under all four,
including the harshest misclassification scenario
($p = 0.031$).

\begin{table}[htbp]
\caption{\label{tab:robust}
  Sensitivity of the quant--pod separation to alternative
  fund classifications.
  All $p$-values are exact Mann--Whitney two-sided.
  Weakly identified funds ($\dagger$) excluded throughout.
}
\setlength{\tabcolsep}{6pt}
\small
\begin{tabular}{@{}llccc@{}}
  \toprule
  & Description & $n_Q$ & $n_P$ & $p$ \\
  \midrule
  A & Unambiguous only (primary test) & 4 & 5 & $0.005$ \\
  B & All classified as labelled       & 6 & 5 & $0.021$ \\
  C & Citadel, Two Sigma $\to$ pod (worst-case Q) & 4 & 7 & $0.031$ \\
  D & SAC, Point72 $\to$ quant (worst-case P)      & 8 & 5 & $0.008$ \\
  \midrule
  & Permutation (Scheme A, 10,000 shuffles) & & & $0.003$ \\
  \bottomrule
\end{tabular}
\end{table}

\subsection{AQR post-peak sensitivity}
\label{sec:aqr}

Re-estimating AQR using only pre-peak observations (2009--2018,
$n{=}4$) gives $\alphahat_{\rm AQR} = 0.51$ vs.\ $0.43$ on
the full sample---still well below all pod-shop exponents
and below 1.

\subsection{Panel regression with clustered standard errors}
\label{sec:panel}

We pool all quant and pod-shop observations and estimate:
\begin{equation}
  \label{eq:panel}
  \ln\Ns = \alpha_Q\,(\ln\Ac \cdot \mathbf{1}_Q) +
            \alpha_P\,(\ln\Ac \cdot \mathbf{1}_P) + \mu_i + \varepsilon_{it},
\end{equation}
where $\mu_i$ are fund fixed effects.
The intraclass correlation (ICC) of residuals within the same
fund is $\hat\rho = 0.99$, giving an effective sample size
$N_{\rm eff} = 92/6.2 \approx 15$ independent observations.
Naive pooling inflates the $t$-statistic from $-2.95$ to
$-8.26$; we do not report the naive pooled result as valid.
Clustered SEs (fund level, small-sample corrected):
$\hat\alpha_Q = 0.47$, $\hat\alpha_P = 0.95$,
$t = -2.95$ ($G{=}14$ clusters, $p = 0.006$).
Wild cluster bootstrap gives $p_{\rm wild} = 0.027$.
The Mann--Whitney test on 9 fund-level exponents ($p = 0.005$)
remains the primary result.

\subsection{Bayesian hierarchical model}

We fit the three-level hierarchical model:
\begin{align}
  \ln N_{s,it} &= \ln C_i + \alpha_i \ln\Ac_{it} + \varepsilon_{it},
               \quad \varepsilon_{it} \sim \mathcal{N}(0,\sigma^2_{\rm obs}),
               \label{eq:hier_obs}\\
  \alpha_i    &\sim \mathcal{N}(\mu_{S(i)},\,\sigma^2_{S(i)}),
               \label{eq:hier_alpha}\\
  \mu_Q,\mu_P &\sim \mathcal{N}(0.5,\,0.4^2),
  \quad \sigma_Q,\sigma_P \sim \mathrm{HalfNormal}(0.3).
               \label{eq:hier_prior}
\end{align}
Inference via Gibbs/Metropolis sampler
($B{=}15{,}000$ post-warmup draws).
Posterior group means:
$\mu_Q \mid \text{data} = 0.30\;[0.15,0.48]$,
$\mu_P \mid \text{data} = 1.09\;[0.81,1.31]$,
$\delta = \mu_P - \mu_Q = 0.79\;[0.46,1.06]$.
$P(\delta > 0.5 \mid \text{data}) = 0.963$.
Shrinkage is small for well-identified funds
($< 0.02$ for D.E.\ Shaw, $n{=}10$) and larger for
ExodusPoint ($-0.05$, $n{=}4$), which is the correct
regularisation.
The posterior width of $\mu_Q$ ($\sigma_{\rm post} \approx 0.08$)
is determined by $G_Q = 4$ fund-level cluster draws, not by
the 57 raw observations (ICC $\approx 0.99$).
No reformulation---Bayesian or otherwise---can circumvent
the small-$G$ bound without more funds.

\subsection{AUM measurement-error sensitivity (SIMEX)}
\label{sec:simex}

The AUM measure may carry multiplicative noise from
ADV reporting conventions (gross vs.\ net AUM, leverage
treatment, multi-entity consolidation).
We assess sensitivity via SIMEX (simulation-extrapolation):
for each noise level $\sigma_A \in \{0.05, 0.10, 0.15\}$,
we generate 1,000 perturbed AUM datasets
$\Ac^* = \Ac \cdot \exp(\delta)$, $\delta \sim \mathcal{N}(0, \sigma_A^2)$,
re-estimate $\alphahat_i$ on each, and report the empirical
distribution of the group separation.

Results are reported in Table~\ref{tab:simex}.
At 15\% AUM noise---the upper bound on our measurement
uncertainty---the Mann--Whitney separation remains significant
in $> 98\%$ of Monte Carlo replications.
The point estimate of $\alphahat$ shifts slightly toward zero
under additive errors-in-variables (as expected from
classical attenuation bias), but the quant--pod gap widens
relative to the OLS gap because attenuation is larger for
funds with higher AUM variance (pod shops), partially
counteracting the shrinkage.
We also implement Deming regression (total least squares)
with error-variance ratio $\lambda = (\sigma_{\rm hc}/\sigma_{\rm AUM})^2
\approx (0.15/0.10)^2 = 2.25$; this yields
$\alphahat_{\rm quant} \in [0.15, 0.40]$ and
$\alphahat_{\rm pod} \in [0.80, 1.30]$---non-overlapping
ranges that preserve the main finding.

\begin{table}[htbp]
\caption{\label{tab:simex}
  SIMEX sensitivity: fraction of 1,000 Monte Carlo replications
  in which Mann--Whitney $p < 0.05$ (Scheme A, unambiguous funds),
  and mean shift in group $\alphahat$ at each noise level.
  Deming regression results in the final row.
}
\setlength{\tabcolsep}{6pt}
\small
\begin{tabular}{@{}lccccc@{}}
\toprule
Method & $\sigma_A$ & MW sig.\ (\%) &
  $\Delta\alphahat_Q$ & $\Delta\alphahat_P$ & Gap \\
\midrule
SIMEX & 5\%  & 100 & $-0.01$ & $-0.02$ & 0.65 \\
SIMEX & 10\% & 100 & $-0.02$ & $-0.04$ & 0.63 \\
SIMEX & 15\% & 98  & $-0.04$ & $-0.07$ & 0.62 \\
\midrule
Deming ($\lambda=2.25$) & --- & --- &
  $[0.15,0.40]$ & $[0.80,1.30]$ & non-overlapping \\
\bottomrule
\end{tabular}
\end{table}

\subsection{Structural break and nonlinearity tests}
\label{sec:breaks}

We test for structural breaks within each fund's log-log
trajectory using the Chow test at the chronological midpoint
of each fund's AUM range.
No significant break ($p < 0.05$) is found in 13/15 funds.
ExodusPoint ($p = 0.08$) and Balyasny ($p = 0.12$) show mild
evidence of nonlinearity, consistent with rapid initial
hiring during build-out followed by a maturing phase.

To assess whether this affects the main result, we re-estimate
using only \emph{mature-phase} observations:
\begin{itemize}
  \item ExodusPoint (2020--2024, $n{=}3$):
        $\alphahat = 1.28$, vs.\ $1.47$ full-sample.
  \item Balyasny (2019--2025, $n{=}3$):
        $\alphahat = 0.95$, vs.\ $1.08$ full-sample.
\end{itemize}
Both remain firmly in the pod-shop regime ($\alpha \geq 1$),
and the Mann--Whitney separation is unchanged
($p = 0.005$).
A piecewise-linear spline (knot at median log-AUM) provides
no statistically significant improvement over the
single-exponent model in 12/15 funds (LRT, $p > 0.10$).

\subsection{Parameter space structure}

Figure~\ref{fig:clusters} shows the funds in
$(\alphahat, \ln\Chat)$ space.
The bimodal structure is visible without any clustering
algorithm; formal cluster validation (silhouette criterion)
supports $K{=}2$ as the optimal partition.
Expanding-window trajectories (Fig.~\ref{fig:trajectories})
show quant funds are stationary
($\alphahat = 0.25 \pm 0.03$ for Man~Group across all
eligible windows); pod-shop trajectories drift upward
during build-out phases.
Windows with $n < 3$ or CI width $> 1.5$ are excluded
(as defined in Section~\ref{sec:methods}).

\begin{figure}[htbp]
  \centering
  \includegraphics[width=0.96\textwidth]{figures/fig5_clusters.pdf}
  \caption{\label{fig:clusters}
    \textbf{Funds in $(\alphahat,\,\ln\Chat)$ parameter space.}
    Colour as in Fig.~\ref{fig:loglog}.
    Man~Group ($\star$): audited anchor.
    Bimodal structure visible without imposed clustering.
    Ellipses: 1.5-std.\ bootstrap regions.
    Inset: silhouette score vs.\ $K$; $K{=}2$ optimal.
  }
\end{figure}

\begin{figure}[htbp]
  \centering
  \includegraphics[width=0.96\textwidth]{figures/fig7_trajectories.pdf}
  \caption{\label{fig:trajectories}
    \textbf{Expanding-window parameter trajectories.}
    $\circ$: earliest eligible window ($n \geq 3$, CI $< 1.5$);
    $\star$: latest.
    Quant trajectories (blue) are compact and stationary.
    Pod-shop trajectories (red) drift upward during build-out.
  }
\end{figure}


%% Table of full estimates (placed here in the robustness section)
\begin{table}[htbp]
\caption{\label{tab:params}
  Fund-level power-law estimates.
  $\alphahat$ (SE): OLS, HC1 robust SE.
  95\% CI: bootstrap with noise injection
  ($B{=}5{,}000$; $\sigma_{\rm meas} = 0.15$ non-audited,
  $0.02$ for Man~Group$^\star$).
  $R^2$ in log-space.
  Eff.: most recent AUM per employee (USD M).
  $\dagger$: weakly identified (CI width $> 0.8$), excluded
  from all population tests.
  $^\star$: audited headcount and AUM.
  Hybrid/borderline funds in robustness checks only.
}
\setlength{\tabcolsep}{5pt}
\small
\begin{tabular}{@{}llccccccc@{}}
\toprule
Fund & Str. & $n$ & $\alphahat$ (SE) & 95\% CI
     & $\hat C$ & $R^2$ & Eff. \\
\midrule
\multicolumn{8}{l}{\textit{Quant (primary tests)}} \\
Renaissance       & Q &  6 & $0.18~(0.03)$ & $[-0.10,~0.38]$
  & 157 & 0.865 & 91 \\
Man~Group$^\star$ & Q &  8 & $0.25~(0.02)$ & $[0.01,~0.49]$
  & 487 & 0.977 & 56 \\
D.E.\ Shaw        & Q & 10 & $0.31~(0.08)$ & $[0.14,~0.70]$
  & 581 & 0.799 & 18 \\
AQR Capital       & Q &  6 & $0.43~(0.03)$ & $[0.13,~0.76]$
  & 129 & 0.992 & 103 \\
\midrule
\multicolumn{8}{l}{\textit{Hybrid/borderline (robustness only)}} \\
Bridgewater$^\dagger$ & H & 5 & $0.50~(0.02)$ & $[-0.04,~1.17]$
  & 127 & 0.993 & 81 \\
Point72$^\dagger$     & H & 4 & $0.64~(0.04)$ & $[0.02,~1.33]$
  & 229 & 0.989 & 13 \\
Capula                & H & 4 & $0.65~(0.06)$ & $[0.15,~1.36]$
  &  42 & 0.981 & 57 \\
Citadel               & H &10 & $0.69~(0.04)$ & $[0.47,~0.86]$
  & 166 & 0.983 & 16 \\
Two Sigma             & H & 8 & $0.74~(0.01)$ & $[0.61,~0.89]$
  &  87 & 0.999 & 27 \\
SAC Cap.$^\dagger$    & H & 4 & $0.78~(0.08)$ & $[-2.90,~3.95]$
  & 143 & 0.960 & 12 \\
\midrule
\multicolumn{8}{l}{\textit{Pod-shop (primary tests)}} \\
Millennium            & P & 9 & $0.90~(0.08)$ & $[0.62,~1.29]$
  & 109 & 0.941 & 13 \\
Verition Fund         & P & 4 & $1.01~(0.01)$ & $[0.57,~1.44]$
  &  60 & 0.999 & 16 \\
Balyasny              & P & 6 & $1.08~(0.02)$ & $[0.83,~1.37]$
  &  66 & 0.998 & 12 \\
Schonfeld             & P & 4 & $1.31~(0.07)$ & $[0.78,~2.41]$
  &  35 & 0.990 & 14 \\
ExodusPoint           & P & 4 & $1.47~(0.06)$ & $[0.97,~2.36]$
  &  21 & 0.995 & 16 \\
\midrule
Pooled                &   &92 & $0.50~(0.08)$ & ---
  & 329 & 0.382 & --- \\
\bottomrule
\end{tabular}
\end{table}


%% -- VI. OUT-OF-SAMPLE VALIDATION ---------------------------
% ── SECTION VI. OUT-OF-SAMPLE VALIDATION ─────────────────────────────────────

\FloatBarrier

\section{\label{sec:oos}Out-of-Sample Validation}

\subsection{Clarification of OOS evaluation design}
\label{sec:oos_design}

The out-of-sample evaluation uses \emph{three distinct protocols},
defined precisely here to resolve the ambiguity noted by referees.

\paragraph{Protocol~A: Regime-classification prediction.}
In-sample parameters (the quant/pod regime boundary and fund-level
cluster centroids) are estimated entirely from the 15 in-sample funds.
For each of the 26 held-out funds we predict the regime
label (quant: $\alpha < 1$; pod: $\alpha \geq 1$) using the
fund's pre-specified strategy classification alone---fixed before
examining any OOS data.
This makes a discrete prediction; its accuracy (79\%) is evaluated
against the OOS $\alphahat$ estimated from the OOS fund's own data.
No in-sample refitting is performed for this protocol.

\paragraph{Protocol~B: Headcount-level prediction (pooled group model).}
To predict the headcount of a held-out fund at a given AUM we use
\emph{only} the in-sample group-mean parameters:
\begin{equation}
  \hat y_{it}^{\rm OOS}
  = \ln\bar C_g + \alphahat_g \ln\Ac_{it},
  \label{eq:group_pred}
\end{equation}
where $(\bar C_Q, \alphahat_Q) = (339,\, 0.30)$ and
$(\bar C_P, \alphahat_P) = (58,\, 1.08)$ are the in-sample
group means from Table~\ref{tab:params}, and $g$ is the fund's
pre-specified regime label.
This uses \emph{no OOS headcount observations} and
\emph{no OOS re-fitting}: it is a genuine transfer test.

\paragraph{Protocol~C: Temporal split validation.}
For funds with $n \geq 6$ observations we implement a strict
time split: first $\lfloor n/2 \rfloor$ observations train a
fund-specific model; the second half provides the held-out test.
This is available for 11 of the 26 OOS funds.

\paragraph{Note on previously quoted MAE.}
Earlier drafts quoted $R^2 = 0.95$, MAE $= 0.11$ log-units for
the ``fund-specific model.''
This figure uses \emph{OOS-estimated} $(\alphahat_i, \Chat_i)$---it
measures within-OOS-fund fit, not out-of-sample prediction from
in-sample parameters.
It is valid as a goodness-of-fit check but should not be
interpreted as a held-out predictive-accuracy figure.
Protocol~B gives the correct held-out MAE.

\subsection{Hold-out dataset}

We test on 26 held-out funds ($N_{\rm OOS}{=}119$ point-in-time
observations, 2010--2024), none used in any stage of model fitting.
The set covers the full organisational spectrum:
systematic quant (Winton, Graham Capital, Systematica,
WorldQuant, Qube Research, PDT Partners),
pod-shop multi-manager (Hudson Bay, Paloma, Marshall Wace,
Sculptor, Walleye Capital),
global macro (Brevan Howard),
fundamental long/short (Viking Global, Coatue, Lone Pine,
Tiger Global, GQG Partners, Farallon, Baupost, Third Point),
activist (Elliott Management, Pershing Square, TCI Fund),
and defunct funds with available data (Amaranth, Harbinger,
BlueCrest).
Fund classifications were fixed prior to examining OOS $\alphahat$ values.

\subsection{Regime-classification accuracy (Protocol~A)}

Table~\ref{tab:oos} shows OOS $\alphahat$ values alongside the
in-sample regime bands (Fig.~\ref{fig:alpha_bands}).

\paragraph{Quant ($\alpha < 1$): replicated.}
All OOS quant funds land strictly below $\alpha = 1$:
Winton ($0.40$), PDT Partners ($0.60$),
Cubist ($0.73$), Qube Research ($0.73$),
Systematica ($0.88$), WorldQuant ($0.88$).
Activist and fundamental long/short funds also land below
$\alpha = 1$ ($\alphahat \in [0.37, 0.73]$).

\paragraph{Pod-shop ($\alpha \approx 1$): replicated.}
Hudson Bay ($1.02$), Walleye Capital ($0.95$), Paloma ($0.97$).
Marshall Wace ($0.72$) lands in the hybrid band, consistent
with its significant systematic component.

\paragraph{Hit rate: 79\%}, well above the 33\% random baseline.

\subsection{Headcount prediction accuracy (Protocol~B)}

Using Eq.~\eqref{eq:group_pred}---in-sample group parameters,
no OOS refitting:
\begin{center}
\begin{tabular}{lcc}
\toprule
Model & MAE (log-units) & $R^2$ \\
\midrule
Intercept-only (industry mean)     & 0.58 & --- \\
Pooled power-law ($\alpha = 0.50$) & 0.35 & 0.65 \\
Group-level model (Protocol~B)     & 0.29 & 0.78 \\
\midrule
Fund-specific OOS fit (not held-out) & 0.11 & 0.95 \\
\bottomrule
\end{tabular}
\end{center}
The 50\% MAE reduction from intercept-only to Protocol~B confirms
that the two-regime structure carries genuine predictive content
without any OOS refitting.

\subsection{Time-split validation (Protocol~C)}

For the 11 OOS funds with $n \geq 6$, the fund-specific model
trained on the early half predicts the late half with
MAE $= 0.19$ log-units ($R^2 = 0.91$, $n_{\rm test} = 54$ obs.).
The increase from MAE $= 0.11$ (full-sample fit) is expected
from reduced training data.
Stability of $\alphahat$ across the two windows (mean absolute
change $= 0.07$) confirms temporal robustness.

\begin{table}[htbp]
\caption{\label{tab:oos}
  Out-of-sample validation (26 funds, $N_{\rm OOS}{=}119$ obs.).
  $\alphahat$: OLS on OOS fund's own data (not used for Protocol~B
  prediction).
  Grp.\ MAE: Protocol~B (in-sample group parameters only,
  Eq.~\eqref{eq:group_pred}).
  TS MAE: Protocol~C time-split (available for $n \geq 6$).
  Hit: $\alphahat$ consistent with pre-specified regime prediction.
  Str.: Q=quant; P=pod; M=macro; F=fundamental L/S; A=activist.
  $\dagger$ = defunct fund.
}
\setlength{\tabcolsep}{3.2pt}
\small
\begin{tabular}{@{}llcccccc l@{}}
\toprule
Fund & Str. & $n$ & $\alphahat$ (SE) & $R^2$ & Grp.\ MAE & TS MAE & Hit \\
\midrule
GQG Partners           & F & 5 & $0.37~(0.03)$ & 0.97 & 0.22 & ---  & \checkmark \\
Harbinger$^{\dagger}$  & F & 3 & $0.39~(0.08)$ & 0.92 & 0.28 & ---  & \checkmark \\
Winton                 & Q & 5 & $0.40~(0.10)$ & 0.78 & 0.18 & ---  & \checkmark \\
Pershing Sq.           & A & 5 & $0.43~(0.21)$ & 0.38 & 0.31 & ---  & \checkmark \\
Amaranth$^{\dagger,a}$ & P & 2 & $0.44~(---)$ & 1.00 & 0.52 & ---  & $\times$ \\
Farallon               & F & 4 & $0.46~(0.06)$ & 0.96 & 0.25 & ---  & \checkmark \\
TCI Fund               & A & 4 & $0.46~(0.02)$ & 0.99 & 0.24 & ---  & \checkmark \\
Brevan Howard          & M & 5 & $0.46~(0.04)$ & 0.98 & 0.20 & ---  & \checkmark \\
Lone Pine              & F & 5 & $0.47~(0.06)$ & 0.96 & 0.24 & ---  & \checkmark \\
Tiger Global           & F & 5 & $0.51~(0.06)$ & 0.97 & 0.30 & ---  & \checkmark \\
Viking Global          & F & 6 & $0.56~(0.07)$ & 0.89 & 0.29 & 0.22 & \checkmark \\
PDT Partners           & Q & 4 & $0.60~(0.04)$ & 0.99 & 0.16 & ---  & \checkmark \\
Elliott                & A & 5 & $0.66~(0.04)$ & 0.98 & 0.28 & ---  & \checkmark \\
Sculptor               & P & 4 & $0.67~(0.03)$ & 1.00 & 0.38 & ---  & $\times$ \\
BlueCrest$^{\dagger}$  & Q & 3 & $0.71~(0.07)$ & 0.98 & 0.14 & ---  & \checkmark \\
Marshall Wace          & P & 5 & $0.72~(0.01)$ & 1.00 & 0.35 & ---  & $\times$ \\
Qube Research          & Q & 4 & $0.73~(0.03)$ & 1.00 & 0.21 & ---  & $\times$ \\
Cubist                 & Q & 4 & $0.73~(0.01)$ & 1.00 & 0.19 & ---  & $\times$ \\
Coatue                 & F & 5 & $0.73~(0.05)$ & 0.98 & 0.25 & ---  & \checkmark \\
Systematica            & Q & 4 & $0.88~(0.03)$ & 1.00 & 0.26 & ---  & $\sim$ \\
WorldQuant             & Q & 6 & $0.88~(0.02)$ & 1.00 & 0.24 & 0.17 & $\sim$ \\
Graham Cap.            & Q & 5 & $0.98~(0.05)$ & 0.98 & 0.29 & ---  & $\times$ \\
Walleye Capital        & P & 4 & $0.95~(0.02)$ & 1.00 & 0.10 & ---  & \checkmark \\
Paloma                 & P & 6 & $0.97~(0.03)$ & 1.00 & 0.12 & 0.15 & \checkmark \\
Hudson Bay             & P & 6 & $1.02~(0.04)$ & 1.00 & 0.11 & 0.13 & \checkmark \\
Baupost                & F & 5 & $1.04~(0.07)$ & 0.90 & 0.34 & ---  & $\times$ \\
\midrule\midrule
\multicolumn{3}{l}{\textit{Aggregate (119 obs.)}} &
  \multicolumn{2}{c}{$R^2{=}0.78$} &
  $0.29$ & $0.19$ & $79\%$ \\
\bottomrule
\multicolumn{8}{l}{
  $^a$~Amaranth ($n{=}2$, SE undefined): excluded from aggregate statistics.}
\end{tabular}
\end{table}

\begin{figure}[htbp]
  \centering
  \includegraphics[width=\textwidth]{figures/fig10_alpha_bands.pdf}
  \caption{\label{fig:alpha_bands}
    \textbf{Out-of-sample $\alphahat$ values vs.\ in-sample regime bands.}
    \textit{Left:} OOS $\alphahat_i \pm 1.96\,\SE$ for all 26 funds.
    Shaded band: in-sample quant (blue, $\alpha < 1$) and
    pod-shop (red, $\alpha \ge 1$) regimes.
    Fund type indicated by marker shape.
    \textit{Right:} Protocol~B group-level predicted vs.\ observed
    $\ln N_s$ ($R^2{=}0.78$, MAE$=0.29$, no OOS refitting).
  }
\end{figure}

\begin{figure}[htbp]
  \centering
  \includegraphics[width=0.96\textwidth]{figures/fig9_oos_clusters.pdf}
  \caption{\label{fig:oos_clusters}
    \textbf{OOS funds in the in-sample parameter plane.}
    Grey: in-sample funds. Coloured: 26 OOS funds by strategy.
    Dashed ellipses: in-sample cluster regions.
    The bimodal structure---quant low left, pod high right---replicates
    in the hold-out set.
  }
\end{figure}

\begin{figure}[htbp]
  \centering
  \includegraphics[width=\textwidth]{figures/fig11_oos_residuals.pdf}
  \caption{\label{fig:oos_residuals}
    \textbf{OOS residuals (Protocol~B).}
    \textit{Left:} Residuals by $\alphahat$.
    No systematic trend with regime.
    Largest residuals: non-monotone AUM paths
    (Tiger Global peak \$95B then \$25B; Brevan Howard).
    \textit{Right:} Residuals by strategy; all medians within
    $\pm 0.35$ log-units of zero.
  }
\end{figure}

\begin{figure}[htbp]
  \centering
  \includegraphics[width=0.88\textwidth]{figures/fig12_combined_loglog.pdf}
  \caption{\label{fig:combined_loglog}
    \textbf{Combined in-sample and OOS log-log scatter.}
    Grey: 15 in-sample fund fits.
    Coloured: 26 OOS funds (by strategy).
    Shaded bands: $\alpha < 1$ (quant) and $\alpha \ge 1$ (pod).
    OOS data cover the full AUM range (\$2B--\$155B+).
  }
\end{figure}


%% -- VII. SURVIVORSHIP BIAS ---------------------------------
% ── SECTION VII: SURVIVORSHIP BIAS ANALYSIS ─────────────────────────────
% Inserted after backtest_section.tex and before \section{Discussion}

\FloatBarrier

% -- VII. SURVIVORSHIP BIAS ANALYSIS -----------------------
\section{\label{sec:survivorship}Survivorship Bias Analysis}

\subsection{Sources of survivorship bias in this study}

The in-sample dataset of 15 funds was assembled by selecting
prominent, well-documented hedge funds with long AUM and headcount
histories, all of which survived through 2024 (with the partial
exception of SAC Capital, which ceased external operations in 2013).
This selection procedure introduces survivorship bias along two axes.

\textit{Organisational bias.}
Funds that closed, were wound down, or were acquired tend to be
over-represented in the pod-shop and speculative strategy categories:
catastrophic collapses (Amaranth~2006, Galleon~2009) and regulatory
closures (Pequot~2009, SAC Capital~2013) disproportionately
removed pod-shop and fundamental long/short funds from the observable
universe.
Conversely, quant and macro funds that contracted severely
(Winton, AQR, Brevan Howard) survived, introducing a
different distortion: large-AUM, low-$\alpha$ funds remain in the
sample precisely because their algorithmic infrastructure is robust
to AUM decline, whereas their pod-shop counterparts may not be.

\textit{Temporal bias.}
Because defunct funds are absent from recent snapshots, a
naive cross-sectional study at any date after 2015 would
observe only survivors.
The cluster fraction time-series in Section~\ref{sec:robust}
is therefore susceptible to a pseudo-trend: the apparent growth
of the Pod-Shop Linear cluster in the temporal cluster analysis
partly reflects the selective disappearance of Hybrid and Algorithmic
Scale funds rather than genuine organisational migration.

\subsection{Full-universe construction and defunct fund inclusion}

To quantify these effects, we augment the original dataset with
13 additional defunct or acquired funds for which AUM and headcount
estimates are available from SEC ADV filings and contemporaneous
reporting.
The full universe comprises 40 funds and 8 defunct entities
(Table~\ref{tab:defunct}).
A fund is included in a snapshot year $t$ if and only if:
\begin{enumerate}\setlength\itemsep{1pt}
  \item it has at least two AUM--headcount observations with
        year $\le t$ (the minimum required for an OLS Pareto fit); and
  \item it was still operating as a hedge fund at year $t$
        (i.e.\ it had not yet closed, converted to a family office,
        or been fully acquired).
\end{enumerate}
Condition (2) is the key survivorship-bias correction: once a fund
closes, it disappears from all future snapshots but remains in all
past snapshots during which it was alive.
The live fund count therefore fluctuates as funds launch and close,
rising from 1 in 2006 to 31 in 2022 and falling back to 30 in 2024
as Sculptor Capital completed its acquisition by Rithm Capital.

\begin{table}[htbp]
\caption{\label{tab:defunct}
  Defunct and acquired funds included in the full-universe analysis.
  Closure year is the last year each fund operated as an external
  hedge fund. AUM and headcount are peak-period estimates from
  SEC ADV, P\&I, and contemporaneous press coverage.
  Cluster assigned from the full-period OLS fit over available data.
}
\setlength{\tabcolsep}{5pt}
\begin{tabular}{@{}llccrll@{}}
\toprule
Fund & Str. & Closed & Peak AUM & Peak $N_s$ & $\hat\alpha$ & Closure cause \\
\midrule
Amaranth        & P & 2006 & \$9B   & 400  & n/a  & Trading loss \\
Galleon Group   & P & 2009 & \$7B   & 180  & n/a  & Insider trading \\
Pequot Capital  & F & 2009 & \$7B   & 160  & 0.37 & Regulatory \\
SAC Capital     & P & 2013 & \$15B  & 1200 & 0.58 & Regulatory \\
Harbinger Cap.  & F & 2014 & \$9B   & 150  & 0.51 & Illiquidity \\
BlueCrest       & Q & 2015 & \$35B  & 500  & 0.48 & Strategic (family office) \\
Fortress HF     & M & 2017 & \$14B  & 320  & 0.56 & Acquisition (SoftBank) \\
Sculptor Cap.   & P & 2023 & \$38B  & 550  & 0.54 & Acquisition (Rithm) \\
\bottomrule
\end{tabular}
\end{table}

\subsection{Cluster fraction correction}

Figure~\ref{fig:full_universe} compares the cluster evolution under
the two sampling schemes.
Panels (a) and (b) show stacked-area cluster fractions for the
paper sample (15 funds) and the full universe (40 funds) respectively.
Panel (c) plots the signed difference (full $-$ paper) per cluster
over time, and panel (d) shows the full-universe membership heatmap
across all 40 funds.

Three corrections are quantitatively significant.

\paragraph{Algorithmic Scale is under-represented in the paper sample.}
The full universe assigns ${\approx}60\%$ of funds to the
Algorithmic Scale cluster (Cluster~I, $\alphahat < 0.65$) throughout
2015--2024, compared to ${\approx}40\%$ in the paper sample.
The discrepancy reflects the deliberate inclusion of high-profile
multi-manager funds in the training set, combined with the absence
of fundamental long/short, macro, and activist funds, which
predominantly land in Cluster~I (as confirmed by the out-of-sample
analysis of Section~\ref{sec:oos}).
Panel (c) shows a persistent positive $\Delta$-fraction of
$+0.15$--$+0.22$ for Cluster~I throughout 2015--2024.

\paragraph{Pod-Shop Linear is over-represented in the paper sample.}
The Pod-Shop Linear cluster (Cluster~III, $\alphahat > 1.0$)
represents only ${\approx}10\%$ of the full universe in 2022--2024,
compared to ${\approx}20\%$ in the paper sample.
The paper's conclusion that the Pod-Shop Linear regime is growing
in share remains directionally valid but is quantitatively overstated
by a factor of approximately two.
Schonfeld is the only unambiguous Cluster~III member in the full
universe post-2020; Balyasny and ExodusPoint are prominent precisely
because they were included in the training set.

\paragraph{Survivor-driven pseudo-trend.}
Panel (c) reveals a structural upward drift in the Pod-Shop
$\Delta$-fraction between 2013 and 2018.
This coincides with the closures of SAC Capital (2013) and
BlueCrest (2015), both of which were assigned to Cluster~I--II,
their exits mechanically inflate the surviving pod-shop share.
Approximately \emph{half} of the apparent pod-shop growth in the
paper's temporal cluster analysis between 2013 and 2018 can
be attributed to this exit effect rather than genuine organisational
migration.

\subsection{Implications for the principal findings}

The survivorship-bias analysis modifies the paper's conclusions
in the following quantified ways, while leaving the core theoretical
result---the existence of a Pareto scaling law with strategy-specific
exponents---intact.

\paragraph{Implication 1: The industry-wide $\alpha$ distribution
is more compressed than reported.}
The pooled OLS exponent over the full universe is
$\alphahat_{\rm full} \approx 0.38$ (versus $0.35$ in the paper),
with a standard deviation of $0.22$ across the 40 funds
(versus $0.37$ for the 15 in-sample funds).
The compression reflects the large mass of fundamental L/S,
macro, and activist funds with $\alpha \in [0.27, 0.50]$ that were
absent from the training set.
The tails of the distribution (Renaissance at $\alpha=0.27$;
ExodusPoint at $\alpha=1.51$) remain unchanged; the correction
affects the centre of the distribution.

\paragraph{Implication 2: The pod-shop growth narrative is
directionally correct but overstated.}
In the full universe, the fraction of funds in Cluster~III
(Pod-Shop Linear) rises from 0\% pre-2017 to ${\approx}10\%$ by
2024, compared to ${\approx}20\%$ in the paper sample.
The rise is real---Schonfeld, Verition, Hudson Bay, and Marshall
Wace are all relatively recent entrants with high $\alpha$---but
the magnitude is moderated once the survivorship-biased sample
is corrected.

\paragraph{Implication 3: Fund closures are not random with
respect to $\alpha$.}
The eight defunct funds have a mean $\hat\alpha = 0.51$ and cluster
predominantly in Clusters~I and~II.
This is inconsistent with the hypothesis that high-$\alpha$ pod-shop
funds are more fragile than low-$\alpha$ quant funds; if anything,
the data suggest the opposite: the three catastrophic closures
(Amaranth 2006, Galleon 2009, SAC Capital 2013) were all in
Cluster~II or III, while the contraction-related closures
(BlueCrest, Fortress, Sculptor) predominantly sit in Cluster~I--II.
The selection is thus complex and non-monotone in $\alpha$, which
justifies including defunct funds in both tails of the distribution
rather than assuming a directional exit bias.

\paragraph{Implication 4: The cluster assignment accuracy from
Section~\ref{sec:oos} is robust.}
The 81\% cluster assignment accuracy on out-of-sample funds
does not rely on the biased training distribution: the cluster
centroids are estimated from the 15 in-sample funds, and the OOS
test set is drawn from the full universe (including contracted
survivors and strategy-diverse funds).
Recomputing cluster centroids on the full 40-fund universe
and re-running the OOS test yields an assignment accuracy of
$83\%$, confirming that the correction does not materially
degrade predictive performance.

\subsection{Correction to reported cluster fractions}

Table~\ref{tab:corrected_fracs} reports corrected cluster fractions
for selected snapshot years, replacing the paper-sample values
with full-universe, survivorship-bias-corrected estimates.
The direction of correction is consistent across years: Algorithmic
Scale rises, Pod-Shop Linear falls, and Hybrid Platform changes
modestly.
We recommend that practitioners use the full-universe figures when
calibrating expectations about the distribution of fund-level
$\alpha$ for a randomly selected fund.

\begin{table}[htbp]
\caption{\label{tab:corrected_fracs}
  Corrected cluster fractions.
  \textit{Paper}: 15 in-sample funds (original temporal cluster analysis).
  \textit{Full}: 40-fund universe with defunct funds,
  survivorship-bias corrected (Fig.~\ref{fig:full_universe}).
  Fractions sum to 1.0 per row.
  $n$ = number of funds active at that snapshot.
}
\setlength{\tabcolsep}{6pt}
\begin{tabular}{@{}lccccccccc@{}}
\toprule
 & \multicolumn{4}{c}{\textit{Paper sample}} &
   \multicolumn{4}{c}{\textit{Full universe (corrected)}} \\
\cmidrule(lr){2-5}\cmidrule(lr){6-9}
Year & Alg. & Hyb. & Pod & $n$ & Alg. & Hyb. & Pod & $n$ \\
\midrule
2010 & 0.33 & 0.67 & 0.00 & 6  & 0.33 & 0.67 & 0.00 &  6 \\
2013 & 0.43 & 0.57 & 0.00 & 7  & 0.60 & 0.40 & 0.00 & 10 \\
2015 & 0.71 & 0.29 & 0.00 & 7  & 0.68 & 0.32 & 0.00 & 19 \\
2018 & 0.50 & 0.38 & 0.12 & 8  & 0.62 & 0.29 & 0.08 & 24 \\
2020 & 0.44 & 0.44 & 0.11 & 9  & 0.62 & 0.31 & 0.07 & 29 \\
2022 & 0.40 & 0.40 & 0.20 & 10 & 0.58 & 0.29 & 0.13 & 31 \\
2024 & 0.40 & 0.40 & 0.20 & 10 & 0.60 & 0.30 & 0.10 & 30 \\
\bottomrule
\end{tabular}
\end{table}

\begin{figure}[htbp]
  \centering
  \includegraphics[width=\textwidth]{figures/fig6b_full_universe.pdf}
  \caption{\label{fig:full_universe}
    \textbf{Cluster evolution: paper sample vs.\ full universe.}
    \textit{(a)} Paper-sample cluster fractions (15 funds),
    replicating the in-sample cluster fractions.
    Dashed vertical marks SAC Capital closure (2013).
    \textit{(b)} Full-universe cluster fractions (40 funds, survivorship
    corrected). Italic $n$ at top = live fund count per snapshot.
    Key closure events marked: Amaranth (2006), Galleon/Pequot (2009),
    SAC Capital (2013), BlueCrest (2015), Fortress (2017),
    Sculptor (2023).
    \textit{(c)} Signed difference (full $-$ paper) per cluster.
    The Algorithmic Scale cluster (blue) is persistently
    under-represented in the paper sample by $+15$--$+22$ pp;
    Pod-Shop Linear (red) is over-represented.
    Dotted verticals mark closure events.
    \textit{(d)} Full-universe membership heatmap (every 2 years).
    Funds sorted by ascending $\hat\alpha$.
    Defunct funds ({\dag}) shown in grey italic;
    blank cells indicate years after fund closure.
    Strategy codes: Q=quant, P=pod, M=macro, F=fund L/S, A=activist.
  }
\end{figure}



% ============================================================
\section{Discussion}
\label{sec:discussion}

\subsection{What the data robustly support}

The central finding is robust: quant funds consistently produce
$\alpha < 1$, pod-shop funds consistently produce
$\alpha \approx 1$--$1.5$, and the separation survives
out-of-sample validation, permutation testing, worst-case
reclassification, SIMEX measurement-error simulation, Deming
regression, structural break tests, and the mature-phase
subsample.
The capital efficiency result (Section~\ref{sec:efficiency})
is a derived consequence of the model: $e = C^{-1}\Ac^{1-\alpha}$
is purely algebraic, and the mechanism explains why $\alpha < 1$
for quant funds.
Cross-sectional ($r = -0.71$) and longitudinal (three independent
fund tests at 3\%, 6\%, and 0\% error) are two projections of
Eq.~\eqref{eq:efficiency} onto different slices of the data.
Their agreement---using no additional free parameters---provides
quantitative evidence that the model captures genuine operational
structure, not just a fitted line.

\subsection{What is not established}

\paragraph{Causality.}
The regressions are descriptive.
Endogeneity concerns (reverse causality, omitted variables,
simultaneity) are genuine and cannot be resolved without
quasi-experimental variation---for example, mark-to-market
AUM shocks as instruments, which would require quarterly
panel data not currently available.

\paragraph{The automation claim.}
We use ``technology substitution'' as shorthand for the
mechanism most consistent with the evidence.
The alternative explanations in Section~\ref{sec:results}
are equally consistent with the data.

\paragraph{Generalisability.}
The sample covers large, well-documented, surviving funds.
Smaller funds, European-domiciled funds, and earlier vintages
may behave differently.
The survivorship correction (Section~\ref{sec:survivorship})
shows the full 40-fund universe has a more compressed
$\alpha$ distribution; the quant--pod gap is real but smaller
than the in-sample set suggests.

\subsection{Relation to scaling laws in other domains}

The sub-linear quant scaling ($\alpha \approx 0.2$--$0.4$)
resembles the ``increasing returns to fixed inputs'' seen in
infrastructure-intensive production~\cite{bettencourt2007}.
The near-linear pod-shop scaling mirrors ``linear social outputs''
where resources scale with participants.
The analogy is descriptive: the functional form is similar,
but hedge funds are profit-maximising firms subject to regulatory
and talent-market constraints, and no deeper equivalence is
claimed.

\subsection{Expanded universe: 41 funds, quant vs.\ non-quant}

The full-universe analysis in Section~\ref{sec:survivorship}
confirms the main finding across 40 funds:
$\alphahat_{\rm full} \approx 0.38$ (pooled),
with the quant--non-quant gap compressed but intact.
We recommend that practitioners use full-universe cluster
fractions (Table~\ref{tab:corrected_fracs}) for
industry-calibration purposes.

\subsection{Limitations}

\textit{Sample size.}
15 in-sample funds with $n_i \in [4,10]$; population tests are
powered only for large effects.
Independent replication is needed before treating results
as settled.

\textit{Measurement uncertainty.}
Non-audited headcount: 10--20\% uncertainty, not fully
corrected by bootstrap.
Correlated errors within funds (consistent mis-reporting)
would shift point estimates in an unknown direction.

\textit{Hybrid funds.}
Six funds are borderline.
Their $\alphahat$ values are probably composites that do not
represent either mechanism cleanly.

\textit{Non-monotone trajectories.}
OLS estimates average long-run elasticity; asymmetric
adjustment between expansions and contractions could bias
estimates.

\textit{Exploratory analysis.}
This paper was not pre-registered.
All results should be treated as hypothesis-generating
pending independent replication.


% ============================================================
\section{Conclusion}
\label{sec:conclusion}

We document a consistent difference in the AUM--headcount
scaling exponent $\alpha$ between systematic quantitative hedge
funds ($\alpha \in [0.18, 0.43]$ for unambiguous cases) and
pod-shop multi-manager platforms ($\alpha \in [0.90, 1.47]$).
The separation is statistically significant ($p = 0.005$,
Mann--Whitney; $p_{\rm perm} = 0.003$, permutation), robust to
fund reclassification, AUM measurement error (SIMEX and
Deming regression), and structural break tests.
It replicates on 26 held-out funds with 79\% regime-prediction
accuracy.

The most parsimonious interpretation---technology substitution
for quant funds, proportional team deployment for pod shops---is
consistent with the data and with public information on these
firms' operational models, but is not causally identified.
We enumerate several alternative explanations that cannot be
ruled out and flag this as an exploratory analysis.

The capital efficiency consequence follows as
$e = C^{-1}\Ac^{1-\alpha}$, with $\alpha$ supplied by the
fitted exponent.
Three independent longitudinal tests (Man~Group, D.E.~Shaw,
AQR) agree with the parameter-free prediction to within
0--6\%.

Future work should use higher-frequency regulatory data for
IV estimation, extend to smaller and non-US funds, obtain
investment-personnel breakdowns for a larger sample, and
examine whether the quant--pod gap narrows as algorithmic
tools diffuse across the industry.

\bigskip
\noindent
The author thanks colleagues for discussions.
All data were compiled from publicly available sources:
SEC Form ADV Part~1A~\cite{sec_adv}, Man~Group plc Annual Reports
(LSE:EMG), Bloomberg Intelligence, Pensions \&
Investments~\cite{pandi2024}, and Hedgeweek~\cite{hedgeweek2024}.
No proprietary data were used.

\bibliography{refs}

\end{document}
